\uuid{api1}
\exo7id{1632}
\titre{exo7 1632}
\auteur{barraud}
\organisation{exo7}
\datecreate{2003-09-01}
\isIndication{false}
\isCorrection{false}
\chapitre{Réduction d'endomorphisme, polynôme annulateur}
\sousChapitre{Diagonalisation}
\cpge{mpsi-mp2i.calcul-matriciel-systemes.e, pcsi.determinants.c, ptsi.determinants.a, mpsi-mp2i.matrices-applications-lineaires.b, pcsi.matrices-applications-lineaires.b, ptsi.matrices-applications-lineaires.b, mpsi-mp2i.determinants.d, pcsi.calcul-matriciel-systemes-lineaires.d, ptsi.calcul-matriciel-systemes-lineaires.d, bcpst1.matrices, mp-mpi.reduction-endomorphismes.c, pc.reduction.b, psi.reduction-endomorphismes.b, pt.reduction-endomorphismes.b, mp-mpi.reduction-endomorphismes.d, pc.reduction.c, psi.reduction-endomorphismes.c, pt.reduction-endomorphismes.c, mp-mpi.reduction-endomorphismes.i, pc.reduction.d, psi.reduction-endomorphismes.d, bcpst2.valeurs-propres-vecteurs-propres.a, bcpst2.valeurs-propres-vecteurs-propres.b}
\module{Algèbre}
\niveau{L2}
\difficulte{}

\contenu{
\question{
Soit $\alpha$ et $\beta$ deux réels, et $A$ la matrice suivante :
$$
  A=\begin{pmatrix}
            1 & -\alpha &  -\alpha &       1 \\
      1-\beta &  \alpha & \alpha-1 &  -\beta \\
        \beta & -\alpha & 1-\alpha & 1+\beta \\
            0 &  \alpha &   \alpha &       0
    \end{pmatrix}
$$
A quelle condition sur $\alpha$ et $\beta$, $A$ est-elle diagonalisable ?

On suppose $\alpha=0$ et $\beta=0$. Vérifier que $A(A-I)=0$. En déduire
$A^{n}$ et $(A+I)^{-1}$ .
}
}
